% problem_id: S001-theorem-smooth-base-change
% source_id: S001 (Stacks Project)
% source_file: etale-cohomology.tex
% statement_locator: etale-cohomology.tex lines 17445--17461
% proof_locator: etale-cohomology.tex lines 17463--17859
% env: theorem
% status: text-extracted-verbatim (difficulty judgement pending, written by hand)
%
% The statement and proof below are the author's own text, extracted verbatim.
% Nothing is paraphrased, shortened, or supplemented.

\begin{theorem}[Smooth base change]
\label{theorem-smooth-base-change}
Consider a cartesian diagram of schemes
$$
\xymatrix{
X \ar[d]_f & Y \ar[l]^h \ar[d]^e \\
S & T \ar[l]_g
}
$$
where $f$ is smooth and $g$ quasi-compact and quasi-separated. Then
$$
f^{-1}R^qg_*\mathcal{F} = R^qh_*e^{-1}\mathcal{F}
$$
for any $q$ and any abelian sheaf $\mathcal{F}$
on $T_\etale$ all of whose stalks at geometric points are torsion of
orders invertible on $S$.
\end{theorem}

\begin{proof}[First proof of smooth base change]
This proof is very long but more direct (using less general theory)
than the second proof given below.

\medskip\noindent
The theorem is local on $X_\etale$. More precisely, suppose we have
$U \to X$ \'etale such that $U \to S$ factors as $U \to V \to S$
with $V \to S$ \'etale. Then we can consider the cartesian square
$$
\xymatrix{
U \ar[d]_{f'} & U \times_X Y \ar[l]^{h'} \ar[d]^{e'} \\
V & V \times_S T \ar[l]_{g'}
}
$$
and setting $\mathcal{F}' = \mathcal{F}|_{V \times_S T}$
we have $f^{-1}R^qg_*\mathcal{F}|_U = (f')^{-1}R^qg'_*\mathcal{F}'$
and $R^qh_*e^{-1}\mathcal{F}|_U = R^qh'_*(e')^{-1}\mathcal{F}'$
(as follows from the compatibility of localization with morphisms of sites, see
Sites, Lemma \ref{sites-lemma-localize-morphism-strong} and
and
Cohomology on Sites, Lemma
\ref{sites-cohomology-lemma-restrict-direct-image-open}).
Thus it suffices to produce an \'etale covering of $X$ by
$U \to X$ and factorizations $U \to V \to S$
as above such that the theorem holds for the diagram with
$f'$, $h'$, $g'$, $e'$.

\medskip\noindent
By the local structure of smooth morphisms, see
Morphisms, Lemma \ref{morphisms-lemma-smooth-etale-over-affine-space},
we may assume $X$ and $S$ are affine and $X \to S$
factors through an \'etale morphism $X \to \mathbf{A}^d_S$.
If we have a tower of cartesian diagrams
$$
\xymatrix{
W \ar[d]_i & Z \ar[l]^j \ar[d]^k \\
X \ar[d]_f & Y \ar[l]^h \ar[d]^e \\
S & T \ar[l]_g
}
$$
and the theorem holds for the bottom and top squares, then
the theorem holds for the outer rectangle; this is formal.
Writing $X \to S$ as the composition
$$
X \to \mathbf{A}^{d - 1}_S \to \mathbf{A}^{d - 2}_S \to \ldots \to
\mathbf{A}^1_S \to S
$$
we conclude that it suffices to prove the theorem when $X$ and $S$
are affine and $X \to S$ has relative dimension $1$.

\medskip\noindent
For every $n \geq 1$ invertible on $S$, let $\mathcal{F}[n]$
be the subsheaf of sections of $\mathcal{F}$ annihilated by $n$. Then
$\mathcal{F} = \colim \mathcal{F}[n]$ by our assumption on
the stalks of $\mathcal{F}$. The functors $e^{-1}$ and $f^{-1}$
commute with colimits as they are left adjoints. The functors
$R^qh_*$ and $R^qg_*$ commute with filtered colimits by
Lemma \ref{lemma-relative-colimit}.
Thus it suffices to prove the theorem for $\mathcal{F}[n]$.
From now on we fix an integer $n$, we work with
sheaves of $\mathbf{Z}/n\mathbf{Z}$-modules and
we assume $S$ is a scheme over $\Spec(\mathbf{Z}[1/n])$.

\medskip\noindent
Next, we reduce to the case where $T$ is affine. Since $g$ is quasi-compact
and quasi-separate and $S$ is affine, the scheme $T$ is quasi-compact and
quasi-separated. Thus we can use the induction principle of
Cohomology of Schemes, Lemma \ref{coherent-lemma-induction-principle}.
Hence it suffices to show that if $T = W \cup W'$
is an open covering and the theorem holds for the squares
$$
\xymatrix{
X \ar[d] & e^{-1}(W) \ar[l]^i \ar[d] \\
S & W \ar[l]_a
}
\quad
\xymatrix{
X \ar[d] & e^{-1}(W') \ar[l]^j \ar[d] \\
S & W' \ar[l]_b
}
\quad
\xymatrix{
X \ar[d] & e^{-1}(W \cap W') \ar[l]^-k \ar[d] \\
S & W \cap W' \ar[l]_c
}
$$
then the theorem holds for the original diagram. To see this we consider the
diagram
$$
\xymatrix{
f^{-1}R^{q - 1}c_*\mathcal{F}|_{W \cap W'} \ar[d]^{\cong} \ar[r] &
f^{-1}R^qg_*\mathcal{F} \ar[d] \ar[r] &
f^{-1}R^qa_*\mathcal{F}|_W \oplus
f^{-1}R^qb_*\mathcal{F}|_{W'} \ar[d]_{\cong} \\
R^qk_*e^{-1}\mathcal{F}|_{e^{-1}(W \cap W')} \ar[r] &
R^qh_*e^{-1}\mathcal{F} \ar[r] &
R^qi_*e^{-1}\mathcal{F}|_{e^{-1}(W)} \oplus
R^qj_*e^{-1}\mathcal{F}|_{e^{-1}(W')}
}
$$
whose rows are the long exact sequences of
Lemma \ref{lemma-relative-mayer-vietoris}.
Thus the $5$-lemma gives the desired conclusion.

\medskip\noindent
Summarizing, we may assume $S$, $X$, $T$, and $Y$ affine,
$\mathcal{F}$ is $n$ torsion, $X \to S$ is smooth of relative dimension $1$,
and $S$ is a scheme over $\mathbf{Z}[1/n]$.
We will prove the theorem by induction on $q$. The base case $q = 0$
is handled by Lemma \ref{lemma-fppf-reduced-fibres-base-change-f-star}.
Assume $q > 0$ and the theorem holds for all smaller degrees.
Choose a short exact sequence
$0 \to \mathcal{F} \to \mathcal{I} \to \mathcal{Q} \to 0$
where $\mathcal{I}$ is an injective sheaf of $\mathbf{Z}/n\mathbf{Z}$-modules.
Consider the induced diagram
$$
\xymatrix{
f^{-1}R^{q - 1}g_*\mathcal{I} \ar[d]_{\cong} \ar[r] &
f^{-1}R^{q - 1}g_*\mathcal{Q} \ar[d]_{\cong} \ar[r] &
f^{-1}R^qg_*\mathcal{F} \ar[d] \ar[r] &
0 \ar[d] \\
R^{q - 1}h_*e^{-1}\mathcal{I} \ar[r] &
R^{q - 1}h_*e^{-1}\mathcal{Q} \ar[r] &
R^qh_*e^{-1}\mathcal{F} \ar[r] &
R^qh_*e^{-1}\mathcal{I}
}
$$
with exact rows. We have the zero in the right upper corner
as $\mathcal{I}$ is injective. The left two vertical arrows are
isomorphisms by induction hypothesis. Thus it suffices to prove
that $R^qh_*e^{-1}\mathcal{I} = 0$.

\medskip\noindent
Write $S = \Spec(A)$ and $T = \Spec(B)$ and say the morphism
$T \to S$ is given by the ring map $A \to B$. We can write
$A \to B = \colim_{i \in I} (A_i \to B_i)$ as a filtered
colimit of maps of rings of finite type over $\mathbf{Z}[1/n]$
(see Algebra, Lemma \ref{algebra-lemma-limit-no-condition}).
For $i \in I$ we set $S_i = \Spec(A_i)$ and $T_i = \Spec(B_i)$.
For $i$ large enough we can find a smooth morphism $X_i \to S_i$
of relative dimension $1$ such that $X = X_i \times_{S_i} S$, see
Limits, Lemmas \ref{limits-lemma-descend-finite-presentation},
\ref{limits-lemma-descend-smooth}, and
\ref{limits-lemma-descend-dimension-d}. Set $Y_i = X_i \times_{S_i} T_i$
to get squares
$$
\xymatrix{
X_i \ar[d]_{f_i} & Y_i \ar[l]^{h_i} \ar[d]^{e_i} \\
S_i & T_i \ar[l]_{g_i}
}
$$
Observe that $\mathcal{I}_i = (T \to T_i)_*\mathcal{I}$
is an injective sheaf of $\mathbf{Z}/n\mathbf{Z}$-modules on
$T_i$, see Cohomology on Sites, Lemma
\ref{sites-cohomology-lemma-pushforward-injective-flat}.
We have $\mathcal{I} = \colim (T \to T_i)^{-1}\mathcal{I}_i$
by Lemma \ref{lemma-linus-hamann}. Pulling back by $e$ we get
$e^{-1}\mathcal{I} = \colim (Y \to Y_i)^{-1}e_i^{-1}\mathcal{I}_i$.
By Lemma \ref{lemma-relative-colimit-general} applied to the system
of morphisms $Y_i \to X_i$ with limit $Y \to X$ we have
$$
R^qh_*e^{-1}\mathcal{I} =
\colim (X \to X_i)^{-1} R^qh_{i, *} e_i^{-1}\mathcal{I}_i
$$
This reduces us to the case where $T$ and $S$ are affine of finite
type over $\mathbf{Z}[1/n]$.

\medskip\noindent
Summarizing, we have an integer $q \geq 1$ such that the theorem holds
in degrees $< q$, the schemes $S$ and $T$ affine of finite type
type over $\mathbf{Z}[1/n]$, we have
$X \to S$ smooth of relative dimension $1$ with $X$ affine, and
$\mathcal{I}$ is an injective sheaf of $\mathbf{Z}/n\mathbf{Z}$-modules
and we have to show that $R^qh_*e^{-1}\mathcal{I} = 0$.
We will do this by induction on $\dim(T)$.

\medskip\noindent
The base case is $T = \emptyset$, i.e., $\dim(T) < 0$.
If you don't like this, you can take as your base case
the case $\dim(T) = 0$. In this case $T \to S$ is finite
(in fact even $T \to \Spec(\mathbf{Z}[1/n])$ is finite as the
target is Jacobson; details omitted), so
$h$ is finite too and hence has
vanishing higher direct images (see references below).

\medskip\noindent
Assume $\dim(T) = d \geq 0$ and we know the result for all situations where $T$
has lower dimension. Pick $U$ affine and \'etale over $X$ and a section $\xi$
of $R^qh_*q^{-1}\mathcal{I}$ over $U$. We have to show
that $\xi$ is zero. Of course, we may replace $X$ by $U$
(and correspondingly $Y$ by $U \times_X Y$)
and assume $\xi \in H^0(X, R^qh_*e^{-1}\mathcal{I})$.
Moreover, since $R^qh_*e^{-1}\mathcal{I}$ is a sheaf,
it suffices to prove that $\xi$ is zero locally on $X$.
Hence we may replace $X$ by the members of an \'etale covering.
In particular, using Lemma \ref{lemma-higher-direct-images}
we may assume that $\xi$ is the image of an element
$\tilde \xi \in H^q(Y, e^{-1}\mathcal{I})$.
In terms of $\tilde \xi$ our task is to show that
$\tilde \xi$ dies in $H^q(U_i \times_X Y, e^{-1}\mathcal{I})$
for some \'etale covering $\{U_i \to X\}$.

\medskip\noindent
By More on Morphisms, Lemma \ref{more-morphisms-lemma-cover-smooth-by-special}
we may assume that $X \to S$ factors as $X \to V \to S$
where $V \to S$ is \'etale and $X \to V$ is a smooth morphism
of affine schemes of relative dimension $1$, has a section, and
has geometrically connected fibres. Observe that
$\dim(V \times_S T) \leq \dim(T) = d$
for example by More on Algebra, Lemma
\ref{more-algebra-lemma-dimension-etale-extension}.
Hence we may then replace $S$ by $V$ and $T$ by $V \times_S T$
(exactly as in the discussion in the first paragraph of the proof).
Thus we may assume $X \to S$ is smooth of relative dimension $1$,
geometrically connected fibres, and has a section $\sigma : S \to X$.

\medskip\noindent
Let $\pi : T' \to T$ be a finite surjective morphism.
We will use below that $\dim(T') \leq \dim(T) = d$, see
Algebra, Lemma \ref{algebra-lemma-integral-dim-up}.
Choose an injective map $\pi^{-1}\mathcal{I} \to \mathcal{I}'$
into an injective sheaf of $\mathbf{Z}/n\mathbf{Z}$-modules.
Then $\mathcal{I} \to \pi_*\mathcal{I}'$ is injective
and hence has a splitting (as $\mathcal{I}$ is an injective
sheaf of $\mathbf{Z}/n\mathbf{Z}$-modules).
Denote $\pi' : Y' = Y \times_T T' \to Y$ the base change of $\pi$
and $e' : Y' \to T'$ the base change of $e$. Picture
$$
\xymatrix{
X \ar[d]_f & Y \ar[l]^h \ar[d]^e & Y' \ar[l]^{\pi'} \ar[d]^{e'} \\
S & T \ar[l]_g & T' \ar[l]_\pi
}
$$
By Proposition \ref{proposition-finite-higher-direct-image-zero} and
Lemma \ref{lemma-finite-pushforward-commutes-with-base-change} we have
$R\pi'_*(e')^{-1}\mathcal{I}' = e^{-1}\pi_*\mathcal{I}'$.
Thus by the Leray spectral sequence
(Cohomology on Sites, Lemma \ref{sites-cohomology-lemma-Leray})
we have
$$
H^q(Y', (e')^{-1}\mathcal{I}') =
H^q(Y, e^{-1}\pi_*\mathcal{I}') \supset H^q(Y, e^{-1}\mathcal{I})
$$
and this remains true after base change by any $U \to X$ \'etale.
Thus we may replace $T$ by $T'$, $\mathcal{I}$ by $\mathcal{I}'$
and $\tilde \xi$ by its image in $H^q(Y', (e')^{-1}\mathcal{I}')$.

\medskip\noindent
Suppose we have a factorization $T \to S' \to S$ where
$\pi : S' \to S$ is finite.
Setting $X' = S' \times_S X$ we can consider the induced diagram
$$
\xymatrix{
X \ar[d]_f & X' \ar[l]^{\pi'} \ar[d]^{f'} & Y \ar[l]^{h'} \ar[d]^e \\
S & S' \ar[l]_\pi & T \ar[l]_g
}
$$
Since $\pi'$ has vanishing higher direct images we see that
$R^qh_*e^{-1}\mathcal{I} = \pi'_*R^qh'_*e^{-1}\mathcal{I}$
by the Leray spectral sequence. Hence
$H^0(X, R^qh_*e^{-1}\mathcal{I}) = H^0(X', R^qh'_*e^{-1}\mathcal{I})$.
Thus $\xi$ is zero if and only if the corresponding section of
$R^qh'_*e^{-1}\mathcal{I}$ is zero\footnote{This
step can also be seen another way. Namely, we have to show that
there is an \'etale covering $\{U_i \to X\}$ such that
$\tilde \xi$ dies in $H^q(U_i \times_X Y, e^{-1}\mathcal{I})$.
However, if we prove there is an \'etale covering
$\{U'_j \to X'\}$ such that $\tilde \xi$ dies in
$H^q(U'_i \times_{X'} Y, e^{-1}\mathcal{I})$, then by property (B)
for $X' \to X$ (Lemma \ref{lemma-finite-B}) there exists an \'etale covering
$\{U_i \to X\}$ such that $U_i \times_X X'$ is a disjoint union
of schemes over $X'$ each of which factors through $U'_j$ for some $j$.
Thus we see that $\tilde \xi$ dies in $H^q(U_i \times_X Y, e^{-1}\mathcal{I})$
as desired.}.
Thus we may replace $S$ by $S'$ and $X$ by $X'$.
Observe that $\sigma : S \to X$ base changes to $\sigma' : S' \to X'$
and hence after this replacement it is still true that $X \to S$
has a section $\sigma$ and geometrically connected fibres.

\medskip\noindent
We will use that $S$ and $T$ are Nagata schemes, see
Algebra, Proposition \ref{algebra-proposition-ubiquity-nagata}
which will guarantee
that various normalizations are finite, see
Morphisms, Lemmas \ref{morphisms-lemma-nagata-normalization-finite} and
\ref{morphisms-lemma-nagata-normalization}.
In particular, we may first replace $T$ by its normalization
and then replace $S$ by the normalization of $S$ in $T$.
Then $T \to S$ is a disjoint union of dominant
morphisms of integral normal schemes, see
Morphisms, Lemma \ref{morphisms-lemma-normal-normalization}.
Clearly we may argue one connected component
at a time, hence we may assume $T \to S$ is a
dominant morphism of integral normal schemes.

\medskip\noindent
Let $s \in S$ and $t \in T$ be the generic points. By
Lemma \ref{lemma-smooth-base-change-fields} there exist finite field
extensions $K/\kappa(t)$ and $k/\kappa(s)$ such that $k$ is contained in $K$
and a finite \'etale Galois covering $Z \to X_k$ with Galois group $G$
of order dividing a power of $n$ split over $\sigma(\Spec(k))$
such that $\tilde \xi$ maps to zero in $H^q(Z_K, e^{-1}\mathcal{I}|_{Z_K})$.
Let $T' \to T$ be the normalization of $T$ in $\Spec(K)$
and let $S' \to S$ be the normalization of $S$ in $\Spec(k)$.
Then we obtain a commutative diagram
$$
\xymatrix{
S' \ar[d] & T' \ar[l] \ar[d] \\
S & T \ar[l]
}
$$
whose vertical arrows are finite. By the arguments given above we
may and do replace $S$ and $T$ by $S'$ and $T'$ (and correspondingly
$X$ by $X \times_S S'$ and $Y$ by $Y \times_T T'$). After this replacement
we conclude we have a finite \'etale Galois covering
$Z \to X_s$ of the generic fibre of $X \to S$
with Galois group $G$ of order dividing a power of $n$
split over $\sigma(s)$ such that $\tilde \xi$ maps to zero in
$H^q(Z_t, (Z_t \to Y)^{-1}e^{-1}\mathcal{I})$.
Here $Z_t = Z \times_S t = Z \times_s t = Z \times_{X_s} Y_t$.
Since $n$ is invertible on $S$,
by Fundamental Groups, Lemma \ref{pione-lemma-extend-covering}
we can find a finite \'etale morphism $U \to X$ whose restriction to $X_s$
is $Z$.

\medskip\noindent
At this point we replace $X$ by $U$ and $Y$ by $U \times_X Y$.
After this replacement it may
no longer be the case that the fibres of $X \to S$ are geometrically
connected (there still is a section but we won't use this), but what
we gain is that after this replacement $\tilde \xi$ maps to zero
in $H^q(Y_t, e^{-1}\mathcal{I})$, i.e., $\tilde \xi$ restricts to
zero on the generic fibre of $Y \to T$.

\medskip\noindent
Recall that $t$ is the spectrum of the function field of $T$, i.e.,
as a scheme $t$ is the limit of the nonempty affine open subschemes of $T$.
By Lemma \ref{lemma-directed-colimit-cohomology} we conclude there exists
a nonempty open subscheme $V \subset T$ such that $\tilde \xi$ maps to zero in
$H^q(Y \times_T V, e^{-1}\mathcal{I}|_{Y \times_T V})$.

\medskip\noindent
Denote $Z = T \setminus V$. Consider the diagram
$$
\xymatrix{
Y \times_T Z \ar[d]_{e_Z} \ar[r]_{i'} &
Y \ar[d]_e &
Y \times_T V \ar[l]^{j'} \ar[d]^{e_V} \\
Z \ar[r]^i &
T &
V \ar[l]_j
}
$$
Choose an injection $i^{-1}\mathcal{I} \to \mathcal{I}'$
into an injective sheaf of $\mathbf{Z}/n\mathbf{Z}$-modules on $Z$.
Looking at stalks we see that the map
$$
\mathcal{I} \to j_*\mathcal{I}|_V \oplus i_*\mathcal{I}'
$$
is injective and hence splits as $\mathcal{I}$ is an injective sheaf
of $\mathbf{Z}/n\mathbf{Z}$-modules. Thus it suffices to show that
$\tilde \xi$ maps to zero in
$$
H^q(Y, e^{-1}j_*\mathcal{I}|_V) \oplus
H^q(Y, e^{-1}i_*\mathcal{I}')
$$
at least after replacing $X$ by the members of an \'etale covering.
Observe that
$$
e^{-1}j_*\mathcal{I}|_V = j'_*e_V^{-1}\mathcal{I}|_V,\quad
e^{-1}i_*\mathcal{I}' = i'_*e_Z^{-1}\mathcal{I}'
$$
By induction hypothesis on $q$ we see that
$$
R^aj'_*e_V^{-1}\mathcal{I}|_V = 0, \quad a = 1, \ldots, q - 1
$$
By the Leray spectral sequence for $j'$ and the vanishing
above it follows that
$$
H^q(Y, j'_*(e_V^{-1}\mathcal{I}|_V))
\longrightarrow
H^q(Y \times_T V, e_V^{-1}\mathcal{I}_V) =
H^q(Y \times_T V, e^{-1}\mathcal{I}|_{Y \times_T V})
$$
is injective. Thus the vanishing of the image of $\tilde \xi$
in the first summand above because we know $\tilde \xi$ vanishes
in $H^q(Y \times_T V, e^{-1}\mathcal{I}|_{Y \times_T V})$.
Since $\dim(Z) < \dim(T) = d$ by induction the image of $\tilde \xi$
in the second summand
$$
H^q(Y, e^{-1}i_*\mathcal{I}') =
H^q(Y, i'_*e_Z^{-1}\mathcal{I}') =
H^q(Y \times_T Z, e_Z^{-1}\mathcal{I}')
$$
dies after replacing $X$ by the members of a suitable \'etale covering.
This finishes the proof of the smooth base change theorem.
\end{proof}
